\documentclass[10pt]{article}
\usepackage[a4paper,margin=1.8cm]{geometry}
\usepackage[T1]{fontenc}
\usepackage{lmodern,microtype,graphicx,booktabs,siunitx}
\usepackage[hidelinks]{hyperref}
\usepackage[font=small,labelfont=bf]{caption}
\usepackage{titlesec}
\titleformat*{\section}{\large\bfseries}
\titlespacing*{\section}{0pt}{0.9em}{0.2em}
\setlength{\textfloatsep}{0.8em}
\setlength{\intextsep}{0.6em}
\setlength{\parskip}{0.4em}
\setlength{\parindent}{0pt}
\graphicspath{{../outputs/}}

\title{\vspace{-1.5em}99\% Paint Coverage Is Not an Even Coat:\\Re-measuring the PaintNet Cuboids}
\author{Julius Wahler}
\date{\today}

\begin{document}
\maketitle
\vspace{-1.5em}

\paragraph{Summary.}
PaintNet and MaskPlanner (Tiboni et al.) report paint coverage (PC) of up to 99\% on the
PaintNet cuboids. We simulated the dataset's expert paths with an open spray-deposition
model and measured what PC leaves out: how even the film is and how much paint it takes.
Under our gun assumptions, the expert paths at the papers' \SI{120}{mm} standoff leave
stripes on every part, with only 20\% of the surface within $\pm20\%$ of the mean
thickness, and still score 100\% PC. A path with \SI{10}{mm} of sideways aim error reaches
99\% PC once it sprays $1.3$--$2.6\times$ the paint (medians over 10 parts). In the other direction, refining the
paths raises the share of surface within tolerance from 20\% to 76\% while PC does not move.
PC is a necessary check, not a measure of coating quality.

\section*{What the papers measure}
PC counts a surface point as covered once its film exceeds a bar, the 10th percentile of
the expert path's film, and averages over the points the expert path covers. Two
consequences follow. The expert path scores 100\% by construction, however uneven its own
film. And paint above the bar is free: doubling the flow can only raise PC, never lower it.
The bar is also relative. A striped expert film sets a low bar and a lenient test.

\section*{What we built}
A small open-source simulator for the 1{,}000 cuboids of the extended PaintNet dataset
(Zenodo 14967945). The spray is a Gaussian cone (\ang{25} half-angle by default) with
inverse-square and incidence-angle falloff and back-face culling. Orientations follow the
authors' loader ($+X$ rotated by \texttt{yzx} Euler angles $B,C,A$), which we verified
against SciPy. On top of the papers' PC we report the thickness coefficient of variation
(CV), the share of surface within $\pm20\%$ of the mean, and a stripe CV along $Z$ on the
side faces, the step-over direction of the raster. A differentiable version of the model
lets us refine paths by gradient descent on per-pose speed, standoff and sideways shift.

\section*{Findings}

\textbf{1. Striped expert paths.} With passes about
\SI{105}{mm} apart, a \ang{25} cone paints a spot of radius \SI{56}{mm} at \SI{120}{mm} and
\SI{93}{mm} at \SI{200}{mm}. Only the wider spot overlaps its neighbours. Over all 1{,}000
cuboids (Table~\ref{tab:stats}), the \SI{120}{mm} films have about 4.5 times the thickness CV of the
\SI{200}{mm} ones, and both score 100\% PC.

\begin{table}[h]
\centering\small
\caption{Expert paths of all 1{,}000 cuboids, \ang{25} cone, medians.}
\label{tab:stats}
\begin{tabular}{lcc}
\toprule
 & \SI{120}{mm} standoff & \SI{200}{mm} standoff \\
\midrule
Thickness CV & 0.50 & 0.11 \\
Stripe CV along $Z$, side faces & 0.42 & 0.06 \\
Surface within $\pm20\%$ of mean & 20\% & 96\% \\
Papers' PC bar, $\times$ mean thickness & 0.34 & 0.88 \\
Papers' PC & 100\% & 100\% \\
\bottomrule
\end{tabular}
\end{table}

\textbf{2. Gun assumptions.} The papers state the
\SI{120}{mm} distance but no cone angle. Across 28 gun settings (cone \ang{15}--\ang{45},
standoff \SIrange{120}{250}{mm}, 8 parts) the surface within tolerance ranges from 8\% to
97\% and the PC bar from 0.11 to 0.90 of the mean thickness. PC of the expert path is
100\% in every cell.

\textbf{3. Aim error.} We added smooth
sideways aim error (across the stroke, correlated over about \SI{200}{mm}) to the expert
paths of 10 parts spread over the dataset, then raised the paint flow until each passed
99\% PC (Figure~\ref{fig:error}). At \SI{10}{mm} RMS the median path needs $1.3\times$ the
paint at \SI{200}{mm} and $2.6\times$ at \SI{120}{mm}. At \SI{20}{mm} it needs $2.0\times$
and $8.1\times$, and one part needs $40\times$. At \SI{200}{mm} the surface within
tolerance falls from 96\% to 82\% at \SI{10}{mm} error and to 60\% at \SI{20}{mm}, while
PC reads 99\% throughout. The spread between parts is small at \SI{200}{mm} and large at
\SI{120}{mm}, where the narrow footprint leaves thin spots that only extra paint lifts.

\begin{figure}[h]
\centering
\includegraphics[width=0.9\linewidth]{challenge.png}
\caption{Expert paths of 10 parts with growing sideways aim error, each with the flow
raised until it passes 99\% PC; median with the range over parts shaded. Left: paint
needed relative to the expert path (log scale). Middle and right: evenness of the
passing film.}
\label{fig:error}
\end{figure}

\textbf{4. Refinement.} Refining 10 parts spread over the
dataset raises the surface within tolerance from 20\% to 76\% at \SI{120}{mm} (using
$0.92\times$ the paint) and from 96\% to 99\% at \SI{200}{mm} ($1.02\times$), with stripe
CV falling by a factor of 2.5 to 3.6 (Figure~\ref{fig:refine}). PC stays at 100\% and
99.8\%. At \SI{120}{mm} most of the gain comes from distance: the refiner backs the gun
off to about \SI{152}{mm} on average.

\begin{figure}[h]
\centering
\includegraphics[width=0.9\linewidth]{refinement.png}
\caption{Expert against refined paths for 10 cuboids at \SI{120}{mm} (red) and
\SI{200}{mm} (blue). The dashed line marks no change.}
\label{fig:refine}
\end{figure}

\section*{What this means}
PC answers whether paint reached the surface, relative to an expert who is right by
definition. It cannot rank two paths that both pass, and it rewards spraying more. For
a painting system the costs sit elsewhere: paint consumed, film outside specification,
and rework. We suggest reporting PC together with three numbers the same simulator already
gives: paint used relative to the expert path, the share of surface within a tolerance
band, and stripe CV along the step-over direction. A bar in absolute film thickness,
rather than a percentile of the expert film, would also make PC comparable across guns
and parts.

\section*{Limits}
The gun model is assumed: the \ang{25} cone sets the stripe amplitude, and a cone of
roughly \ang{35} or wider would close the stripes at \SI{120}{mm}. The paint-for-PC trade
(finding 3) appears at both standoffs, though its size depends on the footprint. No
PaintNet or MaskPlanner predictions were evaluated; the aim error is synthetic. About
0.7\% of recorded poses put the gun inside the part, where it paints nothing; the error
sweep leaves them in place. Occlusion is
back-face culling, which is exact for cuboids only. The simulator has not been compared
with measured films.

\section*{Next steps}
\begin{itemize}\itemsep0pt
\item Evaluate the models' predicted paths: paint used, tolerance share, and their sideways
      error in mm, which places them on Figure~\ref{fig:error}.
\item Fix the gun model with the real fan width and flow at \SI{120}{mm}.
\item Extend the simulator to the
      windows, shelves and containers categories.
\end{itemize}

\end{document}
